
\documentclass{article}
\usepackage{colortbl}
\usepackage{makecell}
\usepackage{multirow}
\usepackage{supertabular}

\begin{document}

\newcounter{utterance}

\centering \large Interaction Transcript for game `codenames', experiment `abstract', episode 1 with reflexion{-}Qwen3.5{-}9B.
\vspace{24pt}

{ \footnotesize  \setcounter{utterance}{1}
\setlength{\tabcolsep}{0pt}
\begin{supertabular}{c@{$\;$}|p{.15\linewidth}@{}p{.15\linewidth}p{.15\linewidth}p{.15\linewidth}p{.15\linewidth}p{.15\linewidth}}
    \# & $\;$A & \multicolumn{4}{c}{Game Master} & $\;\:$B\\
    \hline

    \theutterance \stepcounter{utterance}  
    & & & \multicolumn{2}{p{0.3\linewidth}}{
        \cellcolor[rgb]{0.9,0.9,0.9}{
            \makecell[{{p{\linewidth}}}]{
                \texttt{\tiny{[GM$|$GM]}}
                {'hidden': {'team': ['obliviousness'], 'innocent': ['curiosity', 'panache'], 'opponent': [], 'assassin': []}, 'revealed': {'team': {'team': [], 'innocent': [], 'opponent': [], 'assassin': []}, 'opponent': {'team': [], 'innocent': [], 'opponent': [], 'assassin': []}}}
            }
        }
    }
    & & \\ \\

    \theutterance \stepcounter{utterance}  
    & & \multicolumn{4}{p{0.6\linewidth}}{
        \cellcolor[rgb]{0.9,0.9,0.9}{
            \makecell[{{p{\linewidth}}}]{
                \texttt{\tiny{[P1$\langle$GM]}}
                \texttt{You have attempted this same task before and did not succeed. These are the plans you wrote after those attempts.} \\
\texttt{Trial \#1: New plan: I failed because my clues were too broad, causing my partner to guess distractor or opponent words (like "attunement" and "potentiality") instead of the specific targets. This wasted turns and allowed the opponent to catch up. Next time, I will choose a clue word that is semantically tight to only 2{-}3 remaining target words, explicitly excluding any words that could be confused with distractors or opponent terms. I will prioritize high{-}value targets that share a unique, specific association with the clue to minimize the risk of accidental guesses.} \\
\\ 
\texttt{Trial \#2: New plan: The previous attempt was aborted because the clue "neutrality" was too broad and semantically linked to "renegotiation" (which is an opponent word), causing the partner to guess an invalid target. Next time, I will strictly verify that the chosen clue word has no strong semantic association with any opponent, distractor, or kill words before submission. I will select a highly specific clue that connects only to the remaining team words (e.g., "reinstatement" and "lackluster" via "restoration" or "dimness") and ensure the TARGETS list matches exactly those specific words, avoiding any ambiguity that could lead to an opponent or distractor guess.} \\
\\ 
\texttt{Trial \#3: New plan: I failed because I provided clues that were either too broad ("change") or semantically ambiguous, leading my partner to guess distractor words (like "misapprehension") or opponent words instead of the specific targets. Specifically, "change" is too generic and can easily trigger guesses of "renegotiation" (opponent) or "misapprehension" (distractor) before "obliviousness". Next time, I will choose a highly specific, unique synonym or context that is exclusive to the remaining target word "obliviousness" (e.g., "unawareness" or "blindness") to ensure no other word in the game lists is a plausible guess. I will strictly verify that the clue has zero semantic overlap with the opponent's "idiosyncrasy" or distractors like "curiosity" and "panache" before submitting.} \\
\\ 
\\ 
\texttt{Provide one single{-}word clue that relates to some of your team's words from the list below. You can choose to how many and to which words your clue relates to.} \\
\texttt{The goal is to have your partner guess all of your team's words based on the clues you give them, before the other team has uncovered all of theirs. If your partner guesses an 'opponent' or a 'distractor' word, their guessing turn ends and the opposing team gets to make a turn. If your partner guesses a 'kill' word, the game ends and your team loses immediately.} \\
\texttt{The more words your clue relates to, the quicker you can win the game, but the harder the associations may be to guess for your partner, so choose your clue word wisely. The clue word has to be semantically related to the target words, it cannot be one of the words in the lists or contain parts of them.} \\
\\ 
\texttt{Always give your single{-}word clue and your comma{-}separated list of related target words in the following format and make your answers as short as possible, never include any other text than is required in this form:} \\
\\ 
\texttt{CLUE: <WORD>} \\
\texttt{TARGETS: <TARGETS>} \\
\\ 
\texttt{Your team words are: obliviousness.} \\
\texttt{Your opponent's team words are: .} \\
\texttt{Distractor words are: curiosity, panache.} \\
\texttt{The kill words are: .} \\
            }
        }
    }
    & \\ \\

    \theutterance \stepcounter{utterance}  
    & \multicolumn{4}{p{0.6\linewidth}}{
        \cellcolor[rgb]{0.9,0.9,0.9}{
            \makecell[{{p{\linewidth}}}]{
                \texttt{\tiny{[P1$\rangle$GM]}}
                \texttt{CLUE: unawareness} \\
\texttt{TARGETS: obliviousness} \\
            }
        }
    }
    & & \\ \\

    \theutterance \stepcounter{utterance}  
    & & & \multicolumn{2}{p{0.3\linewidth}}{
        \cellcolor[rgb]{0.9,0.9,0.9}{
            \makecell[{{p{\linewidth}}}]{
                \texttt{\tiny{[GM$|$GM]}}
                \texttt{unawareness} \\
            }
        }
    }
    & & \\ \\

    \theutterance \stepcounter{utterance}  
    & & & \multicolumn{2}{p{0.3\linewidth}}{
        \cellcolor[rgb]{0.9,0.9,0.9}{
            \makecell[{{p{\linewidth}}}]{
                \texttt{\tiny{[GM$|$GM]}}
                ['obliviousness']
            }
        }
    }
    & & \\ \\

    \theutterance \stepcounter{utterance}  
    & & & \multicolumn{2}{p{0.3\linewidth}}{
        \cellcolor[rgb]{0.9,0.9,0.9}{
            \makecell[{{p{\linewidth}}}]{
                \texttt{\tiny{[GM$|$GM]}}
                {'word': 'obliviousness', 'assignment': 'team'}
            }
        }
    }
    & & \\ \\

    \theutterance \stepcounter{utterance}  
    & & & \multicolumn{2}{p{0.3\linewidth}}{
        \cellcolor[rgb]{0.9,0.9,0.9}{
            \makecell[{{p{\linewidth}}}]{
                \texttt{\tiny{[GM$|$GM]}}
                \texttt{opponent has won} \\
            }
        }
    }
    & & \\ \\

    \theutterance \stepcounter{utterance}  
    & \multicolumn{4}{p{0.6\linewidth}}{
        \cellcolor[rgb]{0.95,0.95,0.95}{
            \makecell[{{p{\linewidth}}}]{
                \texttt{\tiny{[P2$\rightarrow$P1]}}
                \texttt{Post{-}episode Reflexion (not sent to GM):} \\
\texttt{New plan: I failed because "unawareness" is a direct synonym of "obliviousness," which violates the rule that the clue cannot be one of the target words or contain parts of them. Additionally, it is too generic, increasing the risk of triggering distractors like "curiosity." Next time, I will select a highly specific, non{-}synonymous concept that implies the state of being oblivious without naming it directly (e.g., "blindspot" or "missed signal") and ensure it shares no semantic overlap with "curiosity" or "panache." I will strictly verify that the clue is not a synonym, substring, or obvious antonym of the target before submission.} \\
            }
        }
    }
    & & \\ \\

\end{supertabular}
}

\end{document}
